\documentclass[dvipsnames]{book}

\usepackage[margin=1cm,a5paper]{geometry}

\usepackage{bold-extra}
\usepackage{polski}
\usepackage{iwona}
\usepackage[OT4]{fontenc}

\usepackage{amsmath,amssymb,amsthm}
\usepackage{pgf,tikz}
\usepackage{hyperref}

\pagestyle{empty}
\setlength{\parindent}{0pt}
\setlength{\parskip}{5pt}
\renewcommand\leq\leqslant
\renewcommand\geq\geqslant

\colorlet{maincolor}{ForestGreen}
\usepackage[most]{tcolorbox}
\usepackage{varwidth}

\newtcolorbox{tresc}[3][]{enhanced, sharp corners, boxrule=0mm, colback=maincolor!10,
colframe=black,toptitle=1mm,fonttitle=\bfseries,
colbacktitle=maincolor!25, coltitle=black,
title={Zadanie #2\hfill{\it Opracowanie: #3}},#1}

\newenvironment{problem}[3]{\newpage\begin{tresc}{#1}{#2}{#3}}{\end{tresc}}

% Dodatkowe makra

\usepackage{multicol}
\usepackage{cancel}
\usepackage{bbm}

\def\ind{\mathbbm{1}}

\def\N{\mathbb{N}}
\def\P{\mathbb{P}}
\def\E{\mathbb{E}}
\def\Z{\mathbb{Z}}
\def\Q{\mathbb{Q}}
\def\R{\mathbb{R}}

\def\nwd{\mathrm{NWD}}

\newtheorem*{lemat}{Lemat}

\usepackage{enumitem}
\setlist[enumerate,1]{label={(\alph*)},topsep=0pt}

% ------------------------------------------------------------------
% Kolorowe boxy: definicje (zielone), twierdzenia/obserwacje/wnioski (niebieskie)
% ------------------------------------------------------------------

\colorlet{defcolor}{ForestGreen}
\colorlet{thmcolor}{RoyalBlue}

\newcounter{defcnt}
\newcounter{thmcnt}

% wspólny styl boxa
\tcbset{notatkibox/.style={enhanced, sharp corners, boxrule=0mm, frame hidden,
  toptitle=0.5mm, bottomtitle=0.5mm, left=2mm, right=2mm, top=1.5mm, bottom=1.5mm,
  fonttitle=\bfseries, coltitle=black, breakable, before skip=6pt, after skip=6pt}}

% \begin{definicja}[nazwa]
\NewDocumentEnvironment{definicja}{o}{%
  \refstepcounter{defcnt}%
  \begin{tcolorbox}[notatkibox, colback=defcolor!8, colframe=defcolor!25,
    colbacktitle=defcolor!25,
    title={Definicja \thedefcnt\IfValueT{#1}{\ (#1)}}]%
}{\end{tcolorbox}}

% generator środowisk niebieskich: #1 = nazwa środowiska, #2 = etykieta
\newcommand{\newthmbox}[2]{%
  \NewDocumentEnvironment{#1}{o}{%
    \refstepcounter{thmcnt}%
    \begin{tcolorbox}[notatkibox, colback=thmcolor!8, colframe=thmcolor!25,
      colbacktitle=thmcolor!25,
      title={#2 \thethmcnt\IfValueT{##1}{\ (##1)}}]%
  }{\end{tcolorbox}}%
}

\newthmbox{twierdzenie}{Twierdzenie}
\newthmbox{obserwacja}{Obserwacja}
\newthmbox{wniosek}{Wniosek}
\newthmbox{stwierdzenie}{Stwierdzenie}

% treść zadania -- żółty box, \begin{zadanie}{1.1}
\colorlet{zadcolor}{Dandelion}

\NewDocumentEnvironment{zadanie}{m}{%
  \begin{tcolorbox}[notatkibox, colback=zadcolor!15, colframe=zadcolor!40,
    colbacktitle=zadcolor!40,
    title={Zadanie #1}]%
}{\end{tcolorbox}}

% przykład -- bez boxa, wyróżniony nagłówkiem
\theoremstyle{definition}
\newtheorem*{przyklad}{Przykład}

\renewcommand{\qedsymbol}{$\square$}


\begin{document}

\begin{center}
  {\Large\bfseries Ćwiczenia 1 --- 2026-10-03}
\end{center}

\section*{Sprawy organizacyjne}

\begin{itemize}[topsep=0pt,itemsep=0pt]
  \item email: \texttt{kacper.kurowski.dokt@pw.edu.pl},
  \item aktywność: $0$--$5$ ptk.,
  \item kartkówki: $0$--$10$ ptk.,
  \item kolokwium: $0$--$30$ ptk.,
  \item egz. pisemny: $0$--$50$ ptk.,
  \item egz. ustny.
\end{itemize}

Aby przystąpić do egz. ustnego trzeba mieć:
\begin{enumerate}[label=\arabic*)]
  \item $\geq 25$ ptk. z egz. pisemnego \textbf{oraz}
  \item $\geq 40$ ptk. z całości cz. ćwiczeniowej.
\end{enumerate}

Dopuszczalna $1$ nieusprawiedliwiona nieobecność oraz $3$ usprawiedliwione
nieobecności.

\hrulefill

\begin{zadanie}{1.1}
Rzucamy wielokrotnie zwyczajną kostką do gry. Które z poniższych ciągów
zmiennych losowych są łańcuchami Markowa? Dla tych, które nimi są, wyznaczyć
macierze prawdopodobieństw przejścia.
\begin{enumerate}
  \item Największa liczba $X_n$ uzyskana w pierwszych $n$ rzutach,
  \item Liczba $N_n$ szóstek uzyskana w pierwszych $n$ rzutach.
\end{enumerate}
\end{zadanie}

\paragraph{(a)} Niech $A_n$ --- liczba uzyskana w $n$-tym rzucie; $A_1, A_2,
\ldots$ są iid o rozkładzie jednostajnym na $S = \{1, \ldots, 6\}$. Przyjmujemy
$X_0 := 1$ (bo $\max$ po zbiorze pustym nie jest określony). Wtedy
\[
  X_n = \max\big(A_1, A_2, \ldots, A_n\big) = \max\big(A_n,\, X_{n-1}\big),
\]
czyli $X_{n+1} = \max(A_{n+1}, X_n)$ --- nowy stan zależy od starego stanu i od
jednego świeżego rzutu. Ustalmy $x_0, \ldots, x_{n+1} \in S$ i oznaczmy
\[
  C := \{X_n = x_n,\ X_{n-1} = x_{n-1},\ \ldots,\ X_0 = x_0\}.
\]
Liczymy:
\begin{align*}
  \P\big(X_{n+1} = x_{n+1} \mid C\big)
  &\overset{(1)}{=} \P\big(\max(A_{n+1},\, X_n) = x_{n+1} \mid C\big) \\[1mm]
  &\overset{(2)}{=} \P\big(\max(A_{n+1},\, x_n) = x_{n+1} \mid C\big) \\[1mm]
  &\overset{(3)}{=} \P\big(\max(A_{n+1},\, x_n) = x_{n+1}\big) \\[1mm]
  &\overset{(3')}{=} \P\big(\max(A_{n+1},\, x_n) = x_{n+1} \mid X_n = x_n\big) \\[1mm]
  &\overset{(2')}{=} \P\big(X_{n+1} = x_{n+1} \mid X_n = x_n\big).
\end{align*}

Uzasadnienia poszczególnych kroków:
\begin{description}[topsep=2pt,itemsep=1pt,leftmargin=10mm]
  \item[(1)] Definicja $X_{n+1}$.

  \item[(2)] \emph{Podstawienie $X_n \to x_n$.} Po podstawieniu zdarzenie pod $\P$ należy już do
  $\sigma(A_{n+1})$ --- nie zależy od historii.

  \item[(3)] \emph{Właściwy krok.} Skoro
  $\{\max(A_{n+1}, x_n) = x_{n+1}\} \in \sigma(A_{n+1})$, a rzuty są niezależne,
  to $A_{n+1}$ jest niezależne od
  $\sigma(A_1, \ldots, A_n) \supseteq \sigma(X_0, \ldots, X_n) \ni C$.
  Warunkowanie przez $C$ niczego więc nie zmienia.

  \item[(3$'$), (2$'$)] Te same dwa kroki wykonane w drugą stronę, dla samego
  warunku $\{X_n = x_n\}$.
\end{description}

Warunek Markowa jest zatem spełniony. Wyznaczmy macierz przejścia. Ustalmy
$X_n = i$; wtedy $X_{n+1} = \max(A_{n+1}, i)$, więc:
\begin{itemize}[topsep=2pt,itemsep=1pt]
  \item $j < i$: niemożliwe, bo $X_n$ jest niemalejący, czyli $p(i,j) = 0$;
  \item $j = i$: zachodzi dokładnie gdy $A_{n+1} \leq i$, czyli
  $p(i,i) = \tfrac{i}{6}$;
  \item $j > i$: zachodzi dokładnie gdy $A_{n+1} = j$, czyli
  $p(i,j) = \tfrac{1}{6}$.
\end{itemize}
Ostatecznie
\[
  p(i,j) =
  \begin{cases}
    \tfrac{1}{6} & \colon j > i, \\[1mm]
    \tfrac{i}{6} & \colon i = j, \\[1mm]
    0            & \colon j < i,
  \end{cases}
  \qquad
  P = \big[p(i,j)\big]_{i,j \in \{1,\ldots,6\}}.
\]
Sprawdzenie stochastyczności: $\sum_{j > i} \tfrac{1}{6} + \tfrac{i}{6}
= \tfrac{6-i}{6} + \tfrac{i}{6} = 1$. \checkmark

\paragraph{(b)} Tutaj schemat jest identyczny --- zmienia się tylko funkcja
wiążąca nowy stan ze starym:
\[
  \begin{cases}
    N_n = N_{n-1} + \ind_{\{A_n = 6\}}, \\
    N_0 = 0.
  \end{cases}
\]
Oznaczmy $D := \{N_{n-1} = i_{n-1},\ \ldots,\ N_0 = i_0\}$. Wtedy
\begin{align*}
  \P\big(N_n = i_n \mid D\big)
  &\overset{(1)}{=} \P\big(N_{n-1} + \ind_{\{A_n = 6\}} = i_n \mid D\big) \\[1mm]
  &\overset{(2)}{=} \P\big(i_{n-1} + \ind_{\{A_n = 6\}} = i_n \mid D\big) \\[1mm]
  &\overset{(3)}{=} \P\big(i_{n-1} + \ind_{\{A_n = 6\}} = i_n\big) \\[1mm]
  &\overset{(3')}{=} \P\big(i_{n-1} + \ind_{\{A_n = 6\}} = i_n \mid N_{n-1} = i_{n-1}\big) \\[1mm]
  &\overset{(2')}{=} \P\big(N_n = i_n \mid N_{n-1} = i_{n-1}\big),
\end{align*}
gdzie uzasadnienia kroków $(1)$--$(3')$ są dokładnie te same co w (a): $(2)$ to
podstawienie $N_{n-1} \to i_{n-1}$, a
$(3)$ korzysta z tego, że $\ind_{\{A_n = 6\}} \in \sigma(A_n)$, zaś $A_n$ jest
niezależne od $\sigma(A_1, \ldots, A_{n-1}) \supseteq \sigma(N_0, \ldots, N_{n-1})$.

Czyli $(N_n)_n$ też jest łańcuchem Markowa. Macierz przejścia: ze stanu $i$
albo wyrzucamy szóstkę (prawdopodobieństwo $\tfrac16$) i przechodzimy do $i+1$,
albo nie (prawdopodobieństwo $\tfrac56$) i zostajemy w $i$:
\[
  p(i,j) =
  \begin{cases}
    \tfrac{1}{6} & \colon j = i+1, \\[1mm]
    \tfrac{5}{6} & \colon i = j, \\[1mm]
    0            & \colon \text{wpp},
  \end{cases}
  \qquad
  P = \big[p(i,j)\big]_{i,j \in \N}.
\]

\begin{zadanie}{1.2}
Niech $(\xi_n)_n$ będzie ciągiem niezależnych zmiennych losowych o jednakowym
rozkładzie zadanym funkcją prawdopodobieństwa
\[
  p(k) = \P(\xi_0 = k), \quad k \in \Z.
\]
Niech $X_0$ będzie zmienną losową przyjmującą wartości w zbiorze
$S = \{1, 2, \ldots, N\}$, niezależną od ciągu $(\xi_n)_n$ oraz niech $f$ będzie
pewną funkcją $f \colon S \times \Z \to S$. Zdefiniujmy ciąg zmiennych losowych
$X_1, X_2, \ldots$:
\[
  X_{n+1} = f(X_n, \xi_n), \quad n = 0, 1, 2, \ldots
\]
Pokazać, że ciąg $(X_n)_n$ jest łańcuchem Markowa.
\end{zadanie}

\begin{proof}
Podstawiamy $X_{n+1} = f(X_n, \xi_n)$:
\begin{align*}
  \P\big(X_{n+1} &= i_{n+1} \mid X_n = i_n, \ldots, X_0 = i_0\big) \\
  &= \P\big(f(X_n, \xi_n) = i_{n+1} \mid X_n = i_n, \ldots, X_0 = i_0\big) \\
  &= \P\big(f(i_n, \xi_n) = i_{n+1} \mid X_n = i_n, \ldots, X_0 = i_0\big) \\
  &= \P\big(f(i_n, \xi_n) = i_{n+1} \mid X_n = i_n\big)
   = \P\big(X_{n+1} = i_{n+1} \mid X_n = i_n\big),
\end{align*}
bo wyrażenie $f(i_n, \xi_n)$ nie mówi nic o $X_{n-1}, X_{n-2}, \ldots, X_0$
(zmienna $\xi_n$ jest niezależna od $X_0$, a więc i od $X_1, \ldots, X_n$),
więc możemy się ich pozbyć. Warunek Markowa jest zatem spełniony.
\end{proof}

\begin{zadanie}{1.3}
Rozważmy łańcuch Markowa o przestrzeni stanów $S = \{0,1,2\}$, jednostajnym
(na $S$) rozkładzie początkowym i macierzy przejścia
\[
  \mathbf{P} =
  \begin{bmatrix}
    0           & \tfrac{1}{2} & \tfrac{1}{2} \\[1mm]
    \tfrac{1}{2} & \tfrac{1}{2} & 0 \\[1mm]
    1           & 0            & 0
  \end{bmatrix}
\]
oraz funkcję $f$ taką, że $f(0) = 0$, a $f(1) = f(2) = 1$. Czy $(f(X_n))_n$ jest
łańcuchem Markowa?
\end{zadanie}

\textbf{Nie.} Kluczowa obserwacja: przejście ze stanu $2$ do $0$ jest
\emph{prawie na pewno} (trzeci wiersz macierzy $\mathbf{P}$ to $[1,0,0]$).

Policzmy dwa prawdopodobieństwa, które gdyby $(f(X_n))_n$ było łańcuchem
Markowa musiałyby być równe.

\paragraph{Pierwsze.} Warunek $f(X_1) = 1$ oznacza $X_1 \in \{1,2\}$, a
$f(X_0) = 1$ oznacza $X_0 \in \{1,2\}$. Gdyby $X_0 = 2$, to $X_1 = 0$ p.n., co
jest sprzeczne z $X_1 \in \{1,2\}$; zatem $X_0 = 1$. Ze stanu $1$ idziemy do
$0$ lub do $1$, więc $X_1 = 1$. Stąd
\begin{align*}
  \P\big(f(X_2) = 0 \mid f(X_1) = 1,\ f(X_0) = 1\big)
  &= \P\big(X_2 = 0 \mid X_1 = 1,\ X_0 = 1\big) \\
  &= \P\big(X_2 = 0 \mid X_1 = 1\big) = \tfrac{1}{2}.
\end{align*}

\paragraph{Drugie.} Teraz $X_0 = 0$ i $X_1 \in \{1,2\}$, przy czym
$\P\big(X_1 \in \{1,2\} \mid X_0 = 0\big) = 1$. Zatem
\begin{align*}
  \P\big(f(X_2) &= 0 \mid f(X_1) = 1,\ f(X_0) = 0\big)
  = \P\big(X_2 = 0 \mid X_1 \in \{1,2\},\ X_0 = 0\big) \\[1mm]
  &= \frac{\begin{aligned}
             &\P\big(X_2 = 0 \mid X_1 = 1\big)\P\big(X_1 = 1 \mid X_0 = 0\big) \\
             + &\P\big(X_2 = 0 \mid X_1 = 2\big)\P\big(X_1 = 2 \mid X_0 = 0\big)
           \end{aligned}}
          {\P\big(X_1 \in \{1,2\} \mid X_0 = 0\big)} \\[1mm]
  &= \tfrac{1}{2} \cdot \tfrac{1}{2} + 1 \cdot \tfrac{1}{2} = \tfrac{3}{4}.
\end{align*}

Ponieważ $\tfrac{1}{2} \neq \tfrac{3}{4}$, warunek Markowa nie jest spełniony
--- $(f(X_n))_n$ \textbf{nie} jest łańcuchem Markowa.

\begin{zadanie}{1.4}
Rozważmy jednorodny łańcuch Markowa o zbiorze stanów $S = \{0,1,2\}$, macierzy
prawdopodobieństw przejścia w jednym kroku
\[
  \mathbf{P} =
  \begin{bmatrix}
    0.1 & 0.2 & 0.7 \\
    0.9 & 0.1 & 0 \\
    0.1 & 0.8 & 0.1
  \end{bmatrix}
\]
oraz wektorze prawdopodobieństw początkowych $(0.3, 0.4, 0.3)$. Obliczyć:
\begin{multicols}{2}
\begin{enumerate}
  \item $\P(X_3 = 1 \mid X_2 = 0)$,
  \item $\P(X_0 = 0, X_1 = 1, X_2 = 2)$,
  \item $\P(X_1 = 1, X_2 = 1 \mid X_0 = 2)$,
  \item $\P(X_1 = 2)$,
  \item $\P(X_0 = 0 \mid X_1 = 2)$,
  \item $\E X_1$.
\end{enumerate}
\end{multicols}
\end{zadanie}

\paragraph{(a)}
\[
  \P\big(X_3 = 1 \mid X_2 = 0\big) = \frac{2}{10}.
\]

\paragraph{(b)}
\begin{align*}
  \P\big(X_0 = 0, X_1 = 1, X_2 = 2\big)
  &= \underbrace{\P\big(X_2 = 2 \mid X_1 = 1\big)}_{= 0} \cdot
     \P\big(X_1 = 1 \mid X_0 = 0\big) \cdot \P\big(X_0 = 0\big) \\
  &= 0.
\end{align*}

\paragraph{(c)} Licznik rozkładamy łańcuchowo, po czym korzystamy z warunku
Markowa (skreślenie warunku $X_0 = 2$) --- czynnik $\P(X_0 = 2)$ się skraca:
\begin{align*}
  \P\big(X_1 = 1, X_2 = 1 \mid X_0 = 2\big)
  &= \frac{\P\big(X_2 = 1,\ X_1 = 1,\ X_0 = 2\big)}{\P\big(X_0 = 2\big)} \\[2mm]
  &\hspace{-16mm} = \frac{\P\big(X_2 = 1 \mid X_1 = 1\big) \cdot
                          \P\big(X_1 = 1 \mid X_0 = 2\big) \cdot
                          \cancel{\P\big(X_0 = 2\big)}}
                         {\cancel{\P\big(X_0 = 2\big)}} \\[2mm]
  &= \P\big(X_2 = 1 \mid X_1 = 1\big) \cdot \P\big(X_1 = 1 \mid X_0 = 2\big) \\[1mm]
  &= \frac{1}{10} \cdot \frac{8}{10} = \frac{8}{100} = \frac{2}{25}.
\end{align*}

\paragraph{(d)}
\[
  \P\big(X_1 = 2\big)
  = \sum_{k=0}^{2} \P\big(X_1 = 2 \mid X_0 = k\big) \cdot \P\big(X_0 = k\big)
  = \frac{24}{100}.
\]

\paragraph{(e)}
\[
  \P\big(X_0 = 0 \mid X_1 = 2\big)
  = \frac{\P\big(X_1 = 2 \mid X_0 = 0\big) \cdot \P\big(X_0 = 0\big)}
         {\P\big(X_1 = 2\big)}
  = \frac{\frac{7}{10} \cdot \frac{3}{10}}{\frac{24}{100}}
  = \frac{\frac{21}{100}}{\frac{24}{100}}
  = \frac{21}{24}.
\]

\paragraph{(f)}
\[
  \E X_1 = \sum_{k=0}^{2} \P\big(X_1 = k\big) \cdot k
  = 0 + 1 \cdot \frac{34}{100} + 2 \cdot \frac{24}{100} = \frac{82}{100}.
\]

\begin{zadanie}{1.5}
Niech $(X_n)_n$ będzie ciągiem niezależnych zmiennych losowych o wartościach w
co najwyżej przeliczalnej przestrzeni stanów $S$. Udowodnić, że $(X_n)_n$ jest
łańcuchem Markowa. Przy jakim założeniu jest on jednorodny?
\end{zadanie}

\begin{proof}
Niech $n \in \N$ oraz $s_0, s_1, \ldots, s_n \in S$. Z niezależności
\begin{align*}
  \P\big(X_n = s_n \mid X_{n-1} &= s_{n-1}, \ldots, X_1 = s_1, X_0 = s_0\big) \\
  &\overset{\text{niezal.}}{=} \P\big(X_n = s_n\big) \\
  &\overset{\text{niezal.}}{=} \P\big(X_n = s_n \mid X_{n-1} = s_{n-1}\big),
\end{align*}
czyli warunek Markowa zachodzi.

Jednorodność oznacza, że
\[
  \P\big(X_n = s_n \mid X_{n-1} = s_{n-1}\big) = \P\big(X_1 = s_n \mid X_0 = s_{n-1}\big),
\]
co --- wobec powyższego --- jest równoważne temu, że
\[
  \P\big(X_n = s_n\big) = \P\big(X_1 = s_n\big) \quad \forall n \in \N,
\]
czyli że wszystkie $X_n$ mają ten sam rozkład. Mamy więc jednorodność
$\iff (X_n)_{n \in \N}$ są iid.
\end{proof}

\begin{zadanie}{1.6}
Pokazać, że jeśli $(X_n)_n$ jest jednorodnym łańcuchem Markowa o macierzy
przejścia $\mathbf{P}$, której każdy wiersz jest taki sam, to:
\begin{enumerate}
  \item $\mathbf{P}(n) = \mathbf{P}$ dla dowolnego $n \in \N$,
  \item $\P(X_n = s) = \P(X_1 = s \mid X_0 = s_0)$ dla dowolnych $s, s_0 \in S$,
  \item $X_0, X_1, X_2, \ldots$ są niezależne.
\end{enumerate}
\end{zadanie}

\paragraph{(a)} Indukcja po $n$.

Dla $n = 1$: $P(1) = P$.

Dla $n = 2$: $P(2) = P^2 = \big[p_{i,j}(2)\big]_{i,j \in S}$, gdzie --- korzystając
z tego, że wszystkie wiersze są równe pierwszemu
\[
  p_{i,j}(2) = \sum_{s \in S} p_{i,s} \, p_{s,j}
             = \sum_{s \in S} p_{1,s} \, p_{1,j}
             = p_{1,j} \cdot \underbrace{\sum_{s \in S} p_{1,s}}_{= 1}
             = p_{i,j},
\]
więc $P(2) = P$.

Krok indukcyjny ($n > 2$): załóżmy, że dla $n \geq 2$ zachodzi $P(n) = P$. Wtedy
\[
  P^{n+1} = \big(P^n\big) \cdot P = P \cdot P = P^2 = P.
\]

\paragraph{(b)}
\begin{align*}
  \P\big(X_n = s\big)
  &= \sum_{s_0 \in S} \P\big(X_n = s \mid X_0 = s_0\big) \cdot \P\big(X_0 = s_0\big) \\
  &\overset{\text{(a)}}{=} \sum_{s_0 \in S} \P\big(X_1 = s \mid X_0 = s_0\big) \cdot \P\big(X_0 = s_0\big)
   = \P\big(X_1 = s \mid X_0 = s_0\big),
\end{align*}
bo prawdopodobieństwo $\P\big(X_1 = s \mid X_0 = s_0\big)$ nie zależy od $s_0$
(wszystkie wiersze $\mathbf{P}$ są takie same), a $\sum_{s_0 \in S} \P(X_0 = s_0) = 1$.

\paragraph{(c)} Z warunku Markowa, jednorodności i punktu (b):
\begin{align*}
  \P\big(X_n = s_n,\ X_{n-1} &= s_{n-1},\ \ldots,\ X_0 = s_0\big) \\
  &= \P\big(X_n = s_n \mid X_{n-1} = s_{n-1}\big) \cdot \ldots \\
  &\qquad \cdot \P\big(X_1 = s_1 \mid X_0 = s_0\big) \cdot \P\big(X_0 = s_0\big) \\
  &= \P\big(X_1 = s_n \mid X_0 = s_{n-1}\big) \cdot \ldots \\
  &\qquad \cdot \P\big(X_1 = s_1 \mid X_0 = s_0\big) \cdot \P\big(X_0 = s_0\big) \\
  &= \P\big(X_n = s_n\big) \cdot \ldots \cdot \P\big(X_1 = s_1\big) \cdot
     \P\big(X_0 = s_0\big),
\end{align*}
czyli $X_0, X_1, X_2, \ldots$ są niezależne.

\begin{zadanie}{1.7}
Pokazać, że produkt dwóch macierzy stochastycznych jest macierzą stochastyczną.
\end{zadanie}

\begin{proof}
Niech $A, B$ --- macierze stochastyczne.

\textbf{Teza:} $P = A \cdot B$ jest macierzą stochastyczną.

Po pierwsze,
\[
  p_{ij} = \sum_{k \in S} a_{ik} b_{kj} \geq 0 \qquad \forall i,j \in S,
\]
jako suma iloczynów liczb nieujemnych. Wtedy $\forall i \in S$:
\[
  \sum_{j \in S} p_{ij}
  = \sum_{j \in S} \sum_{k \in S} a_{ik} b_{kj}
  = \sum_{k \in S} a_{ik} \cdot \underbrace{\sum_{j \in S} b_{kj}}_{= 1}
  = \underbrace{\sum_{k \in S} a_{ik}}_{= 1} = 1.
\]
(W szczególności $p_{ij} \in [0,1]$.)
\end{proof}

\begin{zadanie}{1.8}
Niech $(X_n)_{n \geq 0}$ będzie łańcuchem Markowa na przestrzeni stanów
$S = \{0,1\}$ o macierzy przejścia:
\[
  \mathbf{P} =
  \begin{pmatrix}
    1-a & a \\
    b   & 1-b
  \end{pmatrix},
\]
gdzie $a, b > 0$. Definiujemy nowy proces $(Z_n)_{n \geq 1}$ wzorem
$Z_n = (X_{n-1}, X_n)$, $n \geq 1$. Czy $(Z_n)_{n \geq 1}$ jest łańcuchem
Markowa? Jeśli tak, wyznaczyć macierz przejścia.
\end{zadanie}

\begin{zadanie}{1.9}
Niech $(X_k)_{k \in \N}$ będzie łańcuchem Markowa o przestrzeni stanów
$S = \{0,1,2,3\}$ oraz macierzy przejścia:
\[
  \mathbf{P} =
  \begin{pmatrix}
    1/4 & 0   & 1/2 & 1/4 \\
    0   & 1/5 & 0   & 4/5 \\
    0   & 1   & 0   & 0   \\
    1/3 & 1/3 & 0   & 1/3
  \end{pmatrix}.
\]
Definiujemy nowy proces $Z_n = X_n \ind_{\{X_n \in \{2,3\}\}}$, $n \geq 0$.
\begin{enumerate}
  \item Obliczyć $\P(Z_{n+1} = 2 \mid Z_n = 0, Z_{n-1} = 2)$ i
  $\P(Z_{n+1} = 2 \mid Z_n = 0, Z_{n-1} = 3)$, $n \geq 1$.
  \item Czy $(Z_n)_{n \geq 1}$ jest łańcuchem Markowa?
\end{enumerate}
\end{zadanie}

\begin{zadanie}{1.10}
Niech $(X_n)_{n \geq 1}$ będą wynikami kolejnych rzutów symetryczną kostką do
gry. Rozważmy
\[
  Z_n = X_1 \cdot \ldots \cdot X_n \bmod 5.
\]
Wykazać, że ciąg $(Z_n)_{n \geq 1}$ jest jednorodnym łańcuchem Markowa. Znaleźć
przestrzeń stanów oraz macierz przejścia dla tego łańcucha.
\end{zadanie}

\end{document}
